%% Kommandos fuer Tabellen. Entnommen aus The LateX Companion, tabsatz.ps und diversen Dokus:

%%% ---| Farben fuer Tabellen |-------------------
\IfPackageLoaded{xcolor}{
   \colorlet{tableheadcolor}{gray!30}
   \colorlet{tablesubheadcolor}{gray!20}
   \colorlet{tableblackheadcolor}{black!100}
   \colorlet{tableoddrowcolor}{gray!10.0}
   \colorlet{tableevenrowcolor}{white}
}
%%% ---------------------------------------------



%%% ---|Layout der Tabellen |-------------------


% Groesse der Schrift in Tabellen
\newcommand{\tablefontsize}{\small}
\newcommand{\tableheadfontsize}{\small}
\newcommand{\tablesubheadfontsize}{\small}


% Layout der Tabelle: Ausrichtung, Schrift, Zeilenabstand
\newcommand{\tablestylecommon}{%
  \renewcommand{\arraystretch}{1.3}%  Groessere Abstaende zwischen Zeilen
  %\setlength{\tabcolsep}{0pt}%        Seitliches Padding entfernen
	\small%                             Kleine Schrift
  %\sffamily%                         Serifenlose Schrift
  \centering%                         Tabelle zentrieren
}


% Alternierende Farben:
\newcommand\tablealtcolored{%
	%% Verwendung:
	%% \rowcolors{<starting row index>}{<odd row color>}{<even row color>}
	\rowcolors{1}{tableoddrowcolor}{tableevenrowcolor}%
}


\newcommand{\tablestylealtcolored}{%
	\tablestylecommon%
	\tablealtcolored%
}



%%% -| Neue Spaltendefinitionen 'columntypes' |--
%
% Belegte Spaltentypen:
% l - links
% c - zentriert
% r - rechts
% p,m,b  - oben, mittig, unten
% X - tabularx Auto-Spalte

% um Tabellenspalten mit Flattersatz zu setzen, muss \\ vor
% (z.B.) \raggedright geschuetzt werden:
\newcommand{\PreserveBackslash}[1]{\let\temp=\\#1\let\\=\temp}

%% Forderung des KSP-Verlages: keine Worttrennung in einer Tabelle, daher \raggedright statt \RaggedRight
%\renewcommand\multirowsetup{\RaggedRight}
\renewcommand\multirowsetup{\raggedright}

% Spalten mit Flattersatz und definierte Breite:
% m{} -> mittig
% p{} -> oben
% b{} -> unten
%
%\newcolumntype{L}[1]{>{\hsize=#1\hsize\RaggedRight\arraybackslash}X}%
%\newcolumntype{R}[1]{>{\hsize=#1\hsize\RaggedLeft\arraybackslash}X}%
%\newcolumntype{C}[1]{>{\hsize=#1\hsize\Centering\arraybackslash}X}%
%
%% Forderung des KSP-Verlages: keine Worttrennung in einer Tabelle, daher \raggedright statt \RaggedRight etc.
\newcolumntype{L}[1]{>{\raggedright\arraybackslash}p{#1}} % linksbündig mit Breitenangabe
\newcolumntype{C}[1]{>{\centering\arraybackslash}p{#1}} % zentriert mit Breitenangabe
\newcolumntype{R}[1]{>{\raggedleft\arraybackslash}p{#1}} % rechtsbündig mit Breitenangabe
%
\newcolumntype{M}{>{\begin{minipage}[t]{2cm}\raggedright}c<{\end{minipage}}}
%
\newcolumntype{G}[1]{>{\RaggedLeft\arraybackslash}p{#1}}
\newcolumntype{U}[1]{>{\RaggedRight\arraybackslash}p{#1}}
%\newcolumntype{C}[1]{>{\Centering\arraybackslash}p{#1}}

%\newcolumntype{J}{>{\PreserveBackslash\justifying\hspace{0pt}}X} % Blocksatz mit Silbentrennung

\newcolumntype{v}[1]{>{\PreserveBackslash\RaggedRight\hspace{0pt}}p{#1}}
\newcolumntype{x}{>{\PreserveBackslash\raggedright\hspace{0pt}}X} % linksbündiger Flattersatz ohne Silbentrennung
\newcolumntype{F}{>{\PreserveBackslash\RaggedRight\hspace{0pt}}X} % linksbündiger Flattersatz mit Silbentrennung (KSP-Verlag: unerwünscht!!!)
\newcolumntype{Y}{>{\PreserveBackslash\RaggedLeft\hspace{0pt}}X} % rechtsbündiger Flattersatz mit Silbentrennung
\newcolumntype{y}{>{\PreserveBackslash\raggedleft\hspace{0pt}}X} % rechtsbündiger Flattersatz ohne Silbentrennung

% Tabellenspaltentyp fuer den Kopf: (Farbe + Ausrichtung)
\newcolumntype{H}[1]{>{\columncolor{tableheadcolor}}l}
% % Rechtsbuendig :
% \newcolumntype{R}[1]{>{\PreserveBackslash\RaggedLeft\hspace{0pt}}m{#1}}
% \newcolumntype{S}[1]{>{\PreserveBackslash\RaggedLeft\hspace{0pt}}p{#1}}
% % Zentriert :
% \newcolumntype{Z}[1]{>{\PreserveBackslash\Centering\hspace{0pt}}m{#1}}
% \newcolumntype{A}[1]{>{\PreserveBackslash\Centering\hspace{0pt}}p{#1}}

%%% Spalten fuer Mathematik
% serifenlose Matheschrift
\newcolumntype{s}[1]{>{\DC@{.}{,}{#1}\mathsf\bgroup}l<{\egroup\DV@end}}

% aequivalent aus typokurz (fett+grau+links)
% \newcolumntype{H}{>{\fontseries{b}\selectfont%
%     \columncolor[gray]{.8}[6pt][0pt]}l}
%%% --------------------------------------------


\IfPackageLoaded{booktabs}{
	\setlength\abovetopsep{0pt}
	\setlength\aboverulesep{0pt}
	\setlength\belowrulesep{0pt}
	\setlength\belowbottomsep{0pt}
	\setlength{\extrarowheight}{-0.3ex} % Verkleinerung der Abstandes oberhalb des Zeileninhaltes (zum vertikalen Zentrieren)
}



%%% ---|Listen in Tabellen |--------------------
\newcommand{\removeindentation}{%
	\leftmargini=\labelsep%
	\advance\leftmargini by \labelsep%
}
%
\makeatletter
\newcommand\tableitemize{
	\@minipagetrue%
	\removeindentation
}
\makeatother


%% Aufzählungen in einer Tabelle
\newenvironment{tabitemize}{%
\begin{list}{\textbullet}{%
\setlength\topsep{0pt}%
\setlength\parsep{0pt}%
\setlength\itemsep{0pt}%
\setlength\leftmargin{0em}%
\setlength\leftmargin{1em}%
\setlength\labelwidth{0.5em}%
\setlength\labelsep{0.5em}%
}
}{%
\end{list}
}
%%% --------------------------------------------


% Neue Umgebung fuer Tabellen:

\newenvironment{Tabelle}[2][c]{%
  \begingroup%
  \tablestylecommon%
  \begin{longtable}[#1]{#2}
  }{
  \end{longtable}%
  \endgroup%
}


% Ruecksetzten der Aenderungen
\newcommand\tablerestoresettings{%
  \renewcommand{\arraystretch}{1.0}% Abstaende wieder zuruecksetzen
  \normalsize\rmfamily% Schrift wieder zuruecksetzen
}

% Tabellenkopf: Serifenlos+fett+schraeg+Schriftfarbe
\newcommand\tableheadfont{%
  %\tableheadfontsize%
  %\sffamily%
	\bfseries%
  %\slshape%
  %\color{white}%
}

% Tabellenzwischenüberschrift auf nächsten Seiten
\newcommand\tablesubheadfont{%
	%\tablesubheadfontsize%
  %\sffamily%
	\bfseries%
  %\slshape%
}



\newcommand{\tableend}{\arrayrulecolor{black}\hline}

% Tabellenkopf (1=Spaltentyp, 2=Text)
\newcommand{\tableheadrow}[2]{%
   \multicolumn{1}{#1@{}}{%
     \raisebox{.1mm}{% Ausrichtung der Beschriftung
       #2%
     }\rule{0pt}{4mm}}% unsichtbare Linie, die die Kopfzeile hoeher macht
}


\newcommand{\tablesubhead}[2]{%
  \multicolumn{#1}{>{\columncolor{tablesubheadcolor}}l}{\tablesubheadfont #2}%
}

% Tabellenbody (=Inhalt)
\newcommand\tablebody{%
	\small%
	%\sffamily%
	%\upshape%
}

\newcommand\tableheadshaded{%
	\rowcolor{tableheadcolor}%
}

%%% --------------------------------------------

\newlength{\mylen}
\newlength{\adjusthspace}

\newenvironment{tabularc}[2]
{%
	\setlength\mylen{#2/(#1)-\tabcolsep*2-\arrayrulewidth*(#1+1)/(#1)}%
	%\setlength{\adjusthspace}{((#2-1)/2)*\linewidth}
	%\par\noindent
	%\hspace*{-\the\adjusthspace}
	\begin{tabular}%{#2}%
		{*{#1}{v{\the\mylen}}}%
}
{\end{tabular}\par}
